Condition throughout on \(\mathcal T\), so the trained maps are fixed. By \(\texttt{lem:orbit-score-factorization}\), for each arm \(a\) there are orbit-specific score laws \(F_{a,O}\) for which a calibration orbit--score pair and the target orbit--score pair have laws
\[
\mathsf Q_a(\{O\}\times C)=q_{a,O}F_{a,O}(C),\qquad \mathsf P_a(\{O\}\times C)=p_{a,O}F_{a,O}(C).\tag{1}
\]
The \(M\) calibration pairs are independent and identically distributed and independent of the target pair, armwise.

Suppose first that joint orbit balance holds. Then \(q_{a,O}=p_{a,O}\) for every arm and orbit, so \((1)\) gives \(\mathsf Q_a=\mathsf P_a\). Marginalizing over \(O\) proves equality in law of \(S_{i,a}\) and \(S_{0,a}\). Together with the independence assertion in the factorization lemma, this proves that
\[
S_{1,a},\ldots,S_{M,a},S_{0,a}
\]
are conditionally independent and identically distributed for each fixed arm. This forward implication is the elementary equality of two finite mixtures after orbit-constant conditional score kernels have been established. It neither uses nor implies joint exchangeability across arms: the common target and the shared assignment can couple the armwise collections.

For necessity, suppose that for some arm \(a\) and orbit \(O\),
\[
q_{a,O}\ne p_{a,O}.\tag{2}
\]
Use \(\texttt{def:orbit-separating-law}\), take the covariates constant, and fix \(\widehat m_0=\widehat m_1=0\). The potential outcomes lie in \(\{0,K\}\), so they satisfy the envelope. The indicator of \(O\) is invariant under the declared action because \(O\) is a full \(\mathcal G\)-orbit; the other arm is identically zero. Hence the deterministic prototype is already invariant, and the stated uniform group symmetrization preserves that invariance. The resulting law therefore belongs to \(\mathcal P_{\mathcal G,K}\).

The observable-selector support condition makes the selected observed response equal the selected arm-\(a\) potential outcome. Consequently, under this law,
\[
S_{i,a}=K\mathbf 1\{U_{i,a}\in O\},\qquad S_{0,a}=K\mathbf 1\{T^a\in O\}.\tag{3}
\]
By the definitions of the calibration and target orbit masses, \((3)\) gives
\[
\Pr(S_{i,a}=K\mid\mathcal T)=q_{a,O},\qquad \Pr(S_{0,a}=K\mid\mathcal T)=p_{a,O}.\tag{4}
\]
Equation \((2)\) makes the two score laws in \((4)\) unequal. Therefore, if calibration and target score laws are equal for every bounded invariant law and every fixed measurable trained predictor, every orbit-mass equality must hold. Repeating this argument for both arms proves joint orbit balance. Thus the finite-partition necessity is elementary once the realizable bounded separating witness \((3)\) is available; the matched-randomization content is the realization of that witness through legal observable selectors and one common target.